\documentclass[11pt]{letter}
\usepackage[T1]{fontenc}
\usepackage{lmodern}
\usepackage[margin=1in]{geometry}

\address{School of Electronic and Information Engineering\\
Suzhou University of Science and Technology\\
Suzhou 215009, Jiangsu, China}

\begin{document}

\begin{letter}{The Editor-in-Chief\\Indoor and Built Environment}

\opening{Dear Editor,}

We are pleased to submit our manuscript entitled ``Event-Triggered Predictive Reinforcement Learning with Unsupervised Dynamic Event Gating for Energy Efficient HVAC Control'' for consideration as an original research article in the \textit{Indoor and Built Environment}.

This study addresses energy-efficient operation of the HVAC system in a commercial building, which is a dominant determinant of both the energy performance and the thermal comfort of the indoor environment. Fixed-interval reinforcement learning controllers execute policies uniformly even though building cooling loads evolve at different rates, causing redundant setpoint actuation during stable operation and delayed response to load changes. We propose event-triggered predictive reinforcement learning with unsupervised dynamic event gating (ET-PRL), which converts fixed-step policy execution into state-dependent action updates by treating control triggering as online anomaly detection over a predictive state. A multi-scale fused anomaly score and a local--global two-layer adaptive threshold keep the trigger boundary sensitive to abrupt disturbances while adapting to intraday variation and long-term operating-condition drift. The method is evaluated in a calibrated EnergyPlus model of a commercial building driven by measured chiller operation data. Compared with a fixed-step reinforcement learning controller using the same policy network, ET-PRL reduces policy executions by 53.54\% and daily chiller energy use by 2.13\%, while maintaining chilled-water supply-temperature regulation within the comfort-consistent acceptable range.

We believe this work fits the \textit{Indoor and Built Environment} closely. The journal publishes research on the quality of the indoor and built environment and how it affects the health, performance, efficiency, and comfort of the people living or working there. Our study directly targets the chilled-water supply-temperature control of a commercial-building cooling system, which governs both the energy efficiency of the building and the thermal stability of the indoor environment. The reported daily energy reduction, together with the reduced setpoint actuation that lowers actuator wear and suppresses transient efficiency losses, is a practical contribution to more efficient and more stable indoor environmental control.

The manuscript has six authors, all of whom have read and approved the final version and agree with its submission to the \textit{Indoor and Built Environment}. The CRediT authorship contribution statement is included in the manuscript.

This manuscript is original, has not been published previously, and is not under consideration by any other journal. The authors declare that they have no competing interests. Generative AI-assisted tools were used for language editing and formatting support only; the authors reviewed and edited the content and take full responsibility for the publication. Code and data are available at the repository stated in the manuscript.

Thank you for considering our submission.

\vspace{2\baselineskip}
\noindent Sincerely,\\[2\baselineskip]
Qiming Fu\\
School of Electronic and Information Engineering\\
Suzhou University of Science and Technology\\
Suzhou 215009, Jiangsu, China\\
E-mail: fqm@mail.usts.edu.cn\\[1\baselineskip]
Jianping Chen\\
School of Electronic and Information Engineering\\
Suzhou University of Science and Technology\\
Suzhou 215009, Jiangsu, China\\
E-mail: alan@mail.usts.edu.cn

\end{letter}
\end{document}
